% ══════════════════════════════════════════════
% CHƯƠNG 5 — KẾT LUẬN VÀ HƯỚNG PHÁT TRIỂN
% Người viết: TV5 (Test + Docs)
% ══════════════════════════════════════════════

\chapter{Kết luận và Hướng phát triển}

\section{Kết quả đạt được}\label{Sec5.1}

% [TV5] Viết đoạn giới thiệu ngắn tổng kết quá trình thực hiện,
% sau đó điền \checkmark hoặc để trống tùy từng yêu cầu nhóm đã hoàn thành.

% [Điền đoạn giới thiệu vào đây]
Bảng~\ref{table:ketqua} tổng kết các yêu cầu kỹ thuật đã hoàn thành.

\begin{table}[H]
  \begin{center}
    \begin{tabular}{|c|l|c|}
      \hline
      \textbf{STT} & \textbf{Yêu cầu kỹ thuật} & \textbf{Hoàn thành} \\
      \hline
      1  & Thiết kế ERD và tạo các bảng với Constraints    & % \checkmark hoặc để trống \\
      \hline
      2  & Views                                           & % \checkmark hoặc để trống \\
      \hline
      3  & Functions                                       & % \checkmark hoặc để trống \\
      \hline
      4  & Stored Procedures với Transaction/TRY-CATCH     & % \checkmark hoặc để trống \\
      \hline
      5  & Triggers                                        & % \checkmark hoặc để trống \\
      \hline
      6  & Kiến trúc 3 tầng                                & % \checkmark hoặc để trống \\
      \hline
      7  & Entity Framework Core                           & % \checkmark hoặc để trống \\
      \hline
      8  & ADO.NET Connected + Disconnected                & % \checkmark hoặc để trống \\
      \hline
      9  & LINQ to Objects + LINQ to Entities              & % \checkmark hoặc để trống \\
      \hline
      10 & LINQ to XML — Export và Import                  & % \checkmark hoặc để trống \\
      \hline
      11 & RESTful API với JWT và phân quyền               & % \checkmark hoặc để trống \\
      \hline
      12 & Swagger UI và Postman collection                & % \checkmark hoặc để trống \\
      \hline
    \end{tabular}
  \end{center}
  \caption{Tổng kết các yêu cầu kỹ thuật}\label{table:ketqua}
\end{table}

\section{Đánh giá kết quả kiểm thử}\label{Sec5.2}

% [TV5] Điền số lượng test case và kết quả Pass/Fail từ Postman Collection Runner.

\begin{table}[H]
  \begin{center}
    \begin{tabular}{|l|c|c|c|}
      \hline
      \textbf{Nhóm chức năng} & \textbf{Số test} & \textbf{Pass} & \textbf{Fail} \\
      \hline
      % [Nhóm chức năng 1] & [số] & [số] & [số] \\
      \hline
      % [Nhóm chức năng 2] & [số] & [số] & [số] \\
      \hline
      % [Nhóm chức năng 3] & [số] & [số] & [số] \\
      \hline
      % [Nhóm chức năng 4] & [số] & [số] & [số] \\
      \hline
      % [Nhóm chức năng 5] & [số] & [số] & [số] \\
      \hline
      \textbf{Tổng cộng}      & [số] & [số] & [số] \\
      \hline
    \end{tabular}
  \end{center}
  \caption{Kết quả kiểm thử API bằng Postman}\label{table:test}
\end{table}

\section{Hạn chế của hệ thống}\label{Sec5.3}

% [TV5] Liệt kê các hạn chế thực tế mà nhóm nhận thấy trong quá trình làm.
% Cố gắng trung thực và cụ thể thay vì liệt kê chung chung.

\begin{itemize}
  \item % [Hạn chế 1]
  \item % [Hạn chế 2]
  \item % [Hạn chế 3]
  \item % [Hạn chế 4 — thêm nếu cần]
\end{itemize}

\section{Hướng phát triển}\label{Sec5.4}

% [TV5] Đề xuất các hướng cải tiến/mở rộng trong tương lai.
% Mỗi mục nên có tên ngắn gọn (bold) và mô tả rõ cách thực hiện.

\begin{enumerate}
  \item \textbf{[Hướng 1]}: % [Mô tả cách triển khai]
  \item \textbf{[Hướng 2]}: % [Mô tả cách triển khai]
  \item \textbf{[Hướng 3]}: % [Mô tả cách triển khai]
  \item \textbf{[Hướng 4]}: % [Mô tả cách triển khai — thêm nếu cần]
\end{enumerate}

\section{Bài học kinh nghiệm}\label{Sec5.5}

% [TV5] Mỗi thành viên tự điền 2–3 câu về kinh nghiệm rút ra
% từ phần mình đảm nhận (kỹ thuật, quy trình, teamwork...).

\noindent\textbf{TV1 — [Vai trò]:} % [Điền kinh nghiệm]\\
\noindent\textbf{TV2 — [Vai trò]:} % [Điền kinh nghiệm]\\
\noindent\textbf{TV3 — [Vai trò]:} % [Điền kinh nghiệm]\\
\noindent\textbf{TV4 — [Vai trò]:} % [Điền kinh nghiệm]\\
\noindent\textbf{TV5 — [Vai trò]:} % [Điền kinh nghiệm]
